\chapter{Sensitivity Analysis and Interpreting Computer Output}

An LP is rarely solved once and filed away. Prices shift, capacities are
renegotiated, new products are proposed: and the first question is always
the same: does the old plan survive, and what is the change worth?
Sensitivity analysis answers such what-if
questions \emph{from the solved problem alone}, without starting the simplex
method over. Every serious solver prints the answers in a standard
\emph{sensitivity report}, and reading that
report (the interpretation of a computer solution) is a skill tested in its
own right. This chapter builds the report line by line from the machinery
of Chapters~4--7: the end tableau supplies the numbers, duality explains
what they mean.

\section{Core Rules}

\begin{rulebox}{Setting}
Throughout, the primal is a maximisation in resource form:
\[
  \max\; z=c_{1}x_{1}+\dots+c_{n}x_{n}
  \quad\text{s.t.}\quad
  \sum_{j=1}^{n}a_{ij}x_{j}\le b_{i}\;(i=1,\dots,m),\qquad x_{j}\ge 0 .
\]
It has been solved: the optimal plan is $x^{*}$ with value $z^{*}$, and the
optimal dual solution (shadow prices) is $y^{*}=(y_{1},\dots,y_{m})$.
Sensitivity analysis asks how $x^{*}$ and $z^{*}$ react when a single piece
of data (one $b_{i}$ or one $c_{j}$) is changed. Minimum problems and
$\ge$ rows are handled by measuring every change as an \emph{improvement}
of the objective; the boxes on the classic report layout below show how.
\end{rulebox}

\begin{rulebox}{Shadow price and RHS ranging}
The shadow price $y_{i}$ is the rate of change of the
optimal value with respect to capacity $i$:
\[
  \Delta z^{*} \;=\; y_{i}\,\Delta b_{i},
\]
valid as long as the change leaves the \emph{same set of constraints
binding} (the same optimal basis). The interval of $b_{i}$-values for which
this holds is the \emph{validity range}, reported as an allowable increase
and an allowable decrease around the current $b_{i}$. Inside the range the
\emph{plan} $x^{*}$ moves (the basic variables are recomputed), but the
list of binding constraints (and hence every shadow price) stays the
same. Outside the range the basis changes and the report makes no promise.
For a \emph{nonbinding} constraint: $y_{i}=0$, the allowable decrease
equals the slack, and the allowable increase is $\infty$.
\end{rulebox}

\begin{rulebox}{RHS ranging by hand (two variables)}
If the optimum is the corner where constraints $r$ and $s$ bind, replace
$b_{r}$ by a parameter and re-solve the $2\times 2$ system of the two
binding constraints. The corner coordinates become linear functions of
$b_{r}$. The validity range is the set of $b_{r}$ for which
\begin{itemize}[leftmargin=1.6em,itemsep=0.3ex]
  \item both coordinates stay $\ge 0$, and
  \item every other constraint is still satisfied.
\end{itemize}
On this range $z^{*}(b_{r})$ is linear with slope $y_{r}$.
\end{rulebox}

\begin{rulebox}{Objective-coefficient ranging}
Changing a coefficient $c_{j}$ tilts the iso-profit lines but moves no
constraint: the feasible region is unchanged. The current corner stays
optimal as long as the iso-profit slope stays between the slopes of the two
constraints binding there. For a corner of
$a_{11}x_{1}+a_{12}x_{2}\le b_{1}$ and $a_{21}x_{1}+a_{22}x_{2}\le b_{2}$
(slopes $m_{1}\le m_{2}$, both negative) the condition is
\[
  m_{1}\;\le\;-\frac{c_{1}}{c_{2}}\;\le\;m_{2}.
\]
Fixing one coefficient and solving for the other yields its allowable
increase and decrease. Inside the range the plan $x^{*}$ is unchanged and
\[
  \Delta z^{*}\;=\;x_{j}^{*}\,\Delta c_{j}.
\]
At an endpoint of the range two corners are optimal simultaneously
(multiple optima); beyond it the optimum jumps to a neighbouring corner.
\end{rulebox}

\begin{rulebox}{Reduced cost}
The reduced cost of a variable $x_{j}$ is
\[
  r_{j} \;=\; c_{j}-\sum_{i=1}^{m}a_{ij}\,y_{i}
  \;=\;\text{own margin}\;-\;\text{opportunity cost of the resources used}.
\]
At a maximum every $r_{j}\le 0$; basic (produced) variables have
$r_{j}=0$; a nonbasic variable has $r_{j}\le 0$, and $|r_{j}|$ is exactly
how much its objective coefficient must \emph{improve} before producing it
becomes worthwhile. Equivalently: $c_{j}$ may rise by $|r_{j}|$ before the
current plan changes, while any decrease of $c_{j}$ leaves an unproduced
product unproduced. The same formula prices out a \emph{new} product: it
enters the plan if and only if its margin exceeds
$\sum_{i}a_{ij}\,y_{i}$, the value of the capacity it would consume.
\end{rulebox}

\begin{rulebox}{Reading a solver sensitivity report}
Solvers print two blocks. The \textbf{variable block} has one row per
decision variable:
\begin{itemize}[leftmargin=1.6em,itemsep=0.3ex]
  \item \emph{Final value}: the optimal $x_{j}^{*}$;
  \item \emph{Reduced cost}: $r_{j}$ ($0$ for produced variables);
  \item \emph{Objective coefficient} with \emph{allowable increase /
        decrease}: the range of $c_{j}$ over which the plan $x^{*}$ is
        unchanged.
\end{itemize}
The \textbf{constraint block} has one row per constraint:
\begin{itemize}[leftmargin=1.6em,itemsep=0.3ex]
  \item \emph{Final value}: the left-hand side $\sum_{j}a_{ij}x_{j}^{*}$
        (equal to $b_{i}$ exactly when the constraint binds);
  \item \emph{Shadow price}: $y_{i}$;
  \item \emph{Constraint RHS} with \emph{allowable increase / decrease}:
        the range of $b_{i}$ over which $y_{i}$ stays valid.
\end{itemize}
The entry \texttt{1E+30} is solver notation for $\infty$: no finite limit.
\end{rulebox}

\begin{rulebox}{The classic report layout: dual prices and sign conventions}
Many solvers print a plain-text report instead of a spreadsheet. The
objective counts as row~1, so the constraints are rows $2,3,\dots$
\begin{itemize}[leftmargin=1.6em,itemsep=0.3ex]
  \item \texttt{OBJECTIVE FUNCTION VALUE}: $z^{*}$, or $C^{*}$ for a
        minimum problem.
  \item \texttt{VALUE} and \texttt{REDUCED COST}: $x_{j}^{*}$, and the
        amount by which the objective coefficient of $x_{j}$ would have to
        \emph{improve} (rise in a max problem, fall in a min problem)
        before $x_{j}$ becomes positive. It is printed as a nonnegative
        number and is $0$ for basic variables.
  \item \texttt{SLACK OR SURPLUS}: $b_{i}-a_{i}^{\top}x^{*}$ for a $\le$
        row, $a_{i}^{\top}x^{*}-b_{i}$ for a $\ge$ row; never negative, and
        $0$ exactly when the row binds.
  \item \texttt{DUAL PRICES}: the \emph{improvement} of the objective per
        unit increase of $b_{i}$, that is $\partial z^{*}/\partial b_{i}$
        for a max problem and $-\partial C^{*}/\partial b_{i}$ for a min
        problem.
  \item \texttt{RANGES IN WHICH THE BASIS IS UNCHANGED}: current
        coefficient or right-hand side with allowable increase and
        decrease; \texttt{INFINITY} means no limit.
\end{itemize}
\textbf{Sign rule.} Because improvement is measured the same way for both
senses, the sign of a dual price depends on the row type only: raising
$b_{i}$ loosens a $\le$ row, so its dual price is $\ge 0$; it tightens a
$\ge$ row, so its dual price is $\le 0$; an equality row can have either
sign. Link to Chapter~7: for a max problem the dual price is $y_{i}$; for
a min problem it is $-u_{i}$, where $u_{i}=\partial C^{*}/\partial b_{i}$
is the dual variable of the min problem. Rewriting $\min C$ as
$\max\,(-C)$ leaves every dual price and reduced cost unchanged and only
flips the sign of the objective value.
\end{rulebox}

\begin{rulebox}{Filling the gaps in an incomplete report}
\begin{enumerate}[leftmargin=1.8em,itemsep=0.3ex]
  \item \textbf{Variables.} Reduced cost $>0$: the variable is nonbasic,
        \texttt{VALUE} $=0$. Reduced cost $0$: normally basic (a
        nonbasic variable with reduced cost $0$ signals alternative
        optima). A nonbasic variable also shows a one-sided objective
        range: the improving side equals its reduced cost, the other side
        is \texttt{INFINITY}.
  \item \textbf{Rows.} Slack $0$ or a nonzero dual price: the row binds. A
        nonbinding row has dual price $0$ and a one-sided right-hand-side
        range: \texttt{INFINITY} on its loosening side, and on the other
        side its slack (allowable decrease of a $\le$ row) or surplus
        (allowable increase of a $\ge$ row). The converse can fail, since
        a binding row may also show \texttt{INFINITY} on one side; when
        the slack is blank, count: a nondegenerate optimum has exactly as
        many positive basic variables as there are rows.
  \item \textbf{Primal.} Set the nonbasic variables to $0$, solve the
        binding rows for the basic ones, then compute the missing slacks
        and the objective value.
  \item \textbf{Dual.} Write the dual of the LP as stated (sign rules of
        Chapter~7). Nonbinding rows get dual value $0$; the dual
        constraint of every basic variable becomes an equation. Solve, check
        the signs, and check that the dual slack of each nonbasic variable
        equals its printed reduced cost.
  \item \textbf{Report and certify.} Convert to dual prices ($y_{i}$ for
        max, $-u_{i}$ for min) and confirm $W^{*}=Z^{*}$: a feasible
        primal--dual pair with equal values is optimal.
\end{enumerate}
\end{rulebox}

\begin{rulebox}{Rebuilding the end tableau from a report}
Bring the LP to maximum form with $\le$ rows ($\min C$ becomes
$\max\,(-C)$, each $\ge$ row is multiplied by $-1$), with slack $s_{i}$ in
row~$i$. Then:
\begin{itemize}[leftmargin=1.6em,itemsep=0.3ex]
  \item \textbf{basis:} the variables with positive value and the slacks
        of the rows with positive slack or surplus ($m$ in total; if fewer,
        the optimum is degenerate and zero-valued basic variables complete
        the basis);
  \item \textbf{RHS column:} their values; bottom right: $z^{*}$ (for a
        min problem $-C^{*}$);
  \item \textbf{$z$-row:} under $x_{j}$ its reduced cost; under $s_{i}$
        the dual price of row~$i$ in $\le$ form, which is the printed dual
        price for an original $\le$ row and its negative for an original
        $\ge$ row; $0$ under every basic variable;
  \item \textbf{body:} the unit columns of the basic variables. The
        remaining entries form $B^{-1}A$; the report does not contain
        them, and they can be computed only from the model itself.
\end{itemize}
\end{rulebox}

\begin{rulebox}{What-if questions without re-solving}
Assume resource costs are already counted in the margins $c_{j}$. Then:
\begin{itemize}[leftmargin=1.6em,itemsep=0.3ex]
  \item \textbf{Buy more of resource $i$ at a premium $p$ per unit?}
        Worth it iff $p<y_{i}$; each unit nets $y_{i}-p$; guaranteed only
        up to the allowable increase of $b_{i}$.
  \item \textbf{Sell (rent out) a unit of resource $i$ at price $p$?}
        Worth it iff $p>y_{i}$; guaranteed up to the allowable decrease.
  \item \textbf{Launch a new product with margin $m$ using $a_{i}$ units
        of resource $i$?} It enters the plan iff
        $m>\sum_{i}a_{i}y_{i}$.
  \item \textbf{A margin $c_{j}$ changes within its range?} Plan
        unchanged, $z^{*}$ changes by $x_{j}^{*}\,\Delta c_{j}$.
  \item \textbf{A capacity $b_{i}$ changes within its range?}
        $z^{*}$ changes by $y_{i}\,\Delta b_{i}$; the new plan follows
        from re-solving the binding equations only.
\end{itemize}
\end{rulebox}

\begin{rulebox}{Limits of the report; the 100\%-rule}
Every range in the report covers \emph{one} change at a time, all other
data held fixed. For several \emph{simultaneous} changes (all among the
$c_{j}$, or all among the $b_{i}$) the 100\%-rule gives a quick sufficient check: express each
change as a fraction of its allowable increase or decrease; if the
fractions sum to at most $1$ (i.e.\ $100\%$), the report's conclusions
remain valid. If the sum exceeds $1$ the rule is silent: the basis
\emph{may} have changed, and only re-solving settles it. The rule is
sufficient, not necessary.
\end{rulebox}

\begin{rulebox}{Beyond the allowable range: what can still be said}
Measure the objective as improvement ($z^{*}$ for a max problem, $-C^{*}$
for a min problem). As a function of one right-hand side $b_{i}$, all
other data fixed, the optimal improvement is piecewise linear and
\emph{concave}: on the validity range its slope is the dual price $y_{i}$,
and the slope can only fall as $b_{i}$ grows and only rise as $b_{i}$
falls. So for a change that exceeds the allowable increase $I$ or decrease
$D$:
\begin{itemize}[leftmargin=1.6em,itemsep=0.3ex]
  \item the part of the change inside the range is exact: $+y_{i}I$ for
        the first $I$ units of an increase, $-y_{i}D$ for the first $D$
        units of a decrease;
  \item an increase by $\Delta>I$ improves the objective by \emph{at most}
        $y_{i}\Delta$ in total: every unit beyond the range is worth at
        most $y_{i}$;
  \item a decrease by $\Delta>D$ worsens it by \emph{at least}
        $y_{i}\Delta$: every unit beyond the range costs at least $y_{i}$
        (and the LP may even become infeasible);
  \item loosening a row never hurts, so for a $\le$ row an increase
        beyond $I$ still yields at least the exact gain $y_{i}I$.
\end{itemize}
For an objective coefficient the old plan stays feasible, so
$z^{*}+x_{j}^{*}\Delta c_{j}$ is always a \emph{lower} bound for the new
maximum (and $C^{*}+x_{j}^{*}\Delta c_{j}$ an \emph{upper} bound for a new
minimum), inside the range with equality, outside it as a bound.
\end{rulebox}

\section{Worked Examples}

\begin{worked}{Sensitivity analysis on the running example: how far can $100$ machine hours and the profit $40$ move before the answer changes?}
\wstep{Step 1. Recall what stays fixed.}
From Chapter~5 the optimal basis is $\{x_{1},x_{2}\}$ with
\[
  B=\begin{pmatrix}2&1\\1&1\end{pmatrix},
  \qquad
  B^{-1}=\begin{pmatrix}1&-1\\-1&2\end{pmatrix},
  \qquad
  \mathbf y=(10,20) .
\]
Sensitivity analysis asks: over what range of a parameter does this
\emph{basis} remain optimal? Within that range the answer changes smoothly
and predictably; outside it, a different basis takes over and everything
must be recomputed.

\wstep{Step 2. Right-hand-side ranging for machine hours $b_{1}$.}
The basic variables are $\mathbf x_{B}=B^{-1}\mathbf b$:
\[
  \begin{pmatrix}x_{1}\\x_{2}\end{pmatrix}
  =\begin{pmatrix}1&-1\\-1&2\end{pmatrix}\begin{pmatrix}b_{1}\\80\end{pmatrix}
  =\begin{pmatrix}b_{1}-80\\160-b_{1}\end{pmatrix}.
\]
Feasibility requires both to be $\ge0$:
\[
  b_{1}-80\ge0\ \text{and}\ 160-b_{1}\ge0
  \quad\Longrightarrow\quad
  \boxed{80\le b_{1}\le160 .}
\]
At $b_{1}=100$ this gives $(20,60)$ \checkmark, comfortably inside.

\wstep{Step 3. What the range means.}
\emph{Within} $[80,160]$ the shadow price $10$ is exact: each additional
machine hour is worth $10$ euros. Ten more hours, $b_{1}=110$, gives
$z=2600+10\cdot10=2700$ and the plan $(30,50)$ --- check:
$2(30)+50=110$ \checkmark and $30+50=80$ \checkmark, with
$z=1200+1500=2700$ \checkmark.

\emph{Outside} it the price changes. At $b_{1}=160$, $x_{2}$ hits zero;
beyond that, extra machine hours cannot be used because labour has run out,
and the shadow price drops to $0$.
\wtrap{Shadow prices are local. Applying $10$ euros per hour to a proposal
for $200$ machine hours predicts $z=3600$; the truth is that nothing beyond
$160$ has any value, so $z$ stops at $3200$. Always quote the range with the
price, and never extrapolate past it.}

\wstep{Step 4. Objective-coefficient ranging for $c_{1}=40$.}
Optimality requires the $z$-row entries under the non-basic variables
(the slacks) to stay non-negative. Those entries are
$\mathbf y^{\!\top}=\mathbf c_{B}^{\!\top}B^{-1}$:
\[
  (c_{1},30)\begin{pmatrix}1&-1\\-1&2\end{pmatrix}
  =\bigl(c_{1}-30,\;-c_{1}+60\bigr).
\]
\[
  c_{1}-30\ge0\ \text{and}\ 60-c_{1}\ge0
  \quad\Longrightarrow\quad
  \boxed{30\le c_{1}\le60 .}
\]

\wstep{Step 5. Confirm it geometrically.}
The objective's slope is $-c_{1}/30$. The corner $(20,60)$ stays optimal as
long as that slope lies between the slopes of the two binding constraints,
$-2$ (machine) and $-1$ (labour):
\[
  -2\le-\frac{c_{1}}{30}\le-1
  \quad\Longrightarrow\quad
  30\le c_{1}\le60 . \checkmark
\]
At the endpoints the objective becomes parallel to one of the constraints
and multiple optima appear --- precisely the situation of Chapter~6.
\wtrap{Within this range the optimal \emph{plan} $(20,60)$ does not change
at all; only the optimal \emph{value} does, at $20$ euros per unit of
$c_{1}$. Right-hand-side changes move the plan and keep the price;
objective changes keep the plan and move the value. Confusing which is
which is the chapter's main hazard.}

\wstep{Step 6. Reduced costs, and the new-product question.}
Both $x_{1}$ and $x_{2}$ are basic, so both have reduced cost $0$: they are
already in the plan. A \emph{non-basic} variable's reduced cost says how
much its profit coefficient would have to improve before it enters. For the
hypothetical product~3 of Chapter~5 (one hour of each resource, margin
$35$):
\[
  35-\mathbf y^{\!\top}\mathbf a_{3}=35-(10+20)=5>0 ,
\]
so it enters; at a margin of $28$ it would not, and the break-even margin is
exactly $30$.

\wstep{Step 7. Two parameters at once: the 100\% rule.}
Ranging is valid for \textbf{one change at a time}. For simultaneous
changes, add the fractions of each allowable range used up; if the total is
at most $100\%$, the basis certainly stays optimal. Raising $b_{1}$ from
$100$ to $130$ uses $30/60=50\%$ of its upward slack; a change in $b_{2}$
that uses $40\%$ of its own range brings the total to $90\%$, so the basis
survives. Above $100\%$ the rule is silent --- not violated, merely
uninformative, and re-solving is then the only reliable answer.

\wstep{Step 8. What a solver's sensitivity report contains.}
\begin{center}\footnotesize
\begin{tabular}{ll}
  \hline
  column & meaning\\
  \hline
  final value & the optimal $x_{j}$ or the constraint's usage\\
  reduced cost & $0$ if basic; else the margin shortfall\\
  shadow price & value of one more unit of that right-hand side\\
  allowable increase/decrease & the ranges computed above\\
  \hline
\end{tabular}
\end{center}
Every number in such a report has been derived by hand in this example. The
report is not a black box --- it is this chapter, printed.
\end{worked}

\begin{worked}{Complete an incomplete classic-layout report for a minimum problem, certify it, and rebuild the end tableau}
The LP and the solver's report, with the blanks still to be filled:
\[
  \min\; C=4x_{1}-2x_{2}-5x_{3}
  \quad\text{s.t.}\quad
  \begin{array}{rcrl}
    x_{1}-x_{2} &\ge& -3 & \text{(row 2)}\\
    x_{1}+2x_{2}+2x_{3} &\le& 12 & \text{(row 3)}\\
    2x_{2}-x_{3} &\ge& -3 & \text{(row 4)}\\
    x_{3} &\le& 10 & \text{(row 5)}\\
    x_{1},x_{2},x_{3} &\ge& 0 .
  \end{array}
\]
\begin{center}\footnotesize\ttfamily
\begin{tabular}{lrr}
\multicolumn{3}{l}{OBJECTIVE FUNCTION VALUE\quad 1)\ \underline{\hspace{3em}}}\\[3pt]
VARIABLE & VALUE & REDUCED COST\\
X1 & \underline{\hspace{3em}} & 6.000000\\
X2 & \underline{\hspace{3em}} & 0.000000\\
X3 & \underline{\hspace{3em}} & 0.000000\\[3pt]
ROW & SLACK OR SURPLUS & DUAL PRICES\\
2) & \underline{\hspace{3em}} & \underline{\hspace{3em}}\\
3) & 0.000000 & \underline{\hspace{3em}}\\
4) & 0.000000 & \underline{\hspace{3em}}\\
5) & \underline{\hspace{3em}} & \underline{\hspace{3em}}\\
\end{tabular}

\medskip
\begin{tabular}{lrrr}
\multicolumn{4}{l}{OBJ COEFFICIENT RANGES}\\
VARIABLE & CURRENT COEF & ALLOWABLE INCREASE & ALLOWABLE DECREASE\\
X1 & 4.000000 & INFINITY & 6.000000\\
X2 & $-$2.000000 & 12.000000 & 3.000000\\
X3 & $-$5.000000 & 3.000000 & INFINITY\\[3pt]
\multicolumn{4}{l}{RIGHTHAND SIDE RANGES}\\
ROW & CURRENT RHS & ALLOWABLE INCREASE & ALLOWABLE DECREASE\\
2 & $-$3.000000 & 2.000000 & INFINITY\\
3 & 12.000000 & 12.000000 & 6.000000\\
4 & $-$3.000000 & 6.000000 & 3.000000\\
5 & 10.000000 & INFINITY & 5.000000\\
\end{tabular}
\end{center}

\wstep{Step 1. Can the LP be read as a standard maximum problem?}
Origin test: $0\ge -3$, $0\le 12$, $0\ge -3$, $0\le 10$ all hold. So it
can: maximise $-C=-4x_{1}+2x_{2}+5x_{3}$ subject to $-x_{1}+x_{2}\le 3$,
$x_{1}+2x_{2}+2x_{3}\le 12$, $-2x_{2}+x_{3}\le 3$, $x_{3}\le 10$. We need
this form in Step~10.

\wstep{Step 2. Fix what the columns mean before reading a single number.}
This is a \emph{minimum} problem, so a dual price is the decrease of the
cost per unit increase of the right-hand side, and a reduced cost is how far
a cost coefficient must \emph{fall}. Expected signs: rows 3 and 5 ($\le$)
can only have dual prices $\ge 0$, rows 2 and 4 ($\ge$) only $\le 0$.
\wtrap{A negative dual price in this layout is not an error and not a
``negative shadow price of a resource'': it is the normal sign of a
binding $\ge$ row, in a max problem and in a min problem alike.}

\wstep{Step 3. Which variables are basic?}
X1 has reduced cost $6>0$, so it is nonbasic: $x_{1}^{*}=0$. Its objective
range confirms it: decrease $6$ (the improving direction for a cost),
increase \texttt{INFINITY}. X2 and X3 have reduced cost $0$ and two-sided
ranges: basic.

\wstep{Step 4. Which rows bind?}
Rows 3 and 4 show slack $0$: binding. Rows 2 and 5 are blank, so count:
four rows mean four basic variables, $x_{2}$ and $x_{3}$ are two of them,
and (the optimum being nondegenerate) the other two are the slacks of
rows 2 and 5, which are therefore nonbinding. Their ranges fit this:
\texttt{INFINITY} on the loosening side, and on the other side the room
left. Row~2 ($\ge$) has allowable increase $2$, so its surplus should be
$2$; row~5 ($\le$) has allowable decrease $5$, so its slack should be $5$.
Step~5 confirms both.

\wstep{Step 5. Solve the primal.}
With $x_{1}=0$ the binding rows read
\[
  2x_{2}+2x_{3}=12,\qquad 2x_{2}-x_{3}=-3 .
\]
Subtracting: $3x_{3}=15$, so $x_{3}=5$ and $x_{2}=1$. Then
\[
  \text{row 2: } 0-1-(-3)=2\ \text{(surplus)},\qquad
  \text{row 5: } 10-5=5\ \text{(slack)},\qquad
  C^{*}=0-2-25=-27 .
\]
Both values agree with the ranges. \checkmark

\wstep{Step 6. Solve the dual by complementary slackness.}
By the sign rules for a min problem, the $\ge$ rows get $u_{1},u_{3}\ge 0$
and the $\le$ rows get $u_{2},u_{4}\le 0$:
\[
  \max\; W=-3u_{1}+12u_{2}-3u_{3}+10u_{4}
  \quad\text{s.t.}\quad
  \begin{array}{rcl}
    u_{1}+u_{2} &\le& 4\\
    -u_{1}+2u_{2}+2u_{3} &\le& -2\\
    2u_{2}-u_{3}+u_{4} &\le& -5 .
  \end{array}
\]
Rows 2 and 5 are slack, so $u_{1}=u_{4}=0$. $x_{2},x_{3}>0$ turn the
second and third dual constraints into equations:
\[
  2u_{2}+2u_{3}=-2,\qquad 2u_{2}-u_{3}=-5
  \quad\Longrightarrow\quad 3u_{3}=3,\; u_{3}=1,\; u_{2}=-2 .
\]
Signs: $u_{2}=-2\le 0$, $u_{3}=1\ge 0$. \checkmark\ The first dual
constraint: $u_{1}+u_{2}=-2\le 4$, slack $6$, which is exactly the printed
reduced cost of X1. \checkmark

\wstep{Step 7. Fill in the report.}
Dual price $=-u_{i}$ for a min problem:
\begin{center}\footnotesize\ttfamily
\begin{tabular}{lrr}
\multicolumn{3}{l}{OBJECTIVE FUNCTION VALUE\quad 1)\ $-$27.00000}\\[3pt]
VARIABLE & VALUE & REDUCED COST\\
X1 & 0.000000 & 6.000000\\
X2 & 1.000000 & 0.000000\\
X3 & 5.000000 & 0.000000\\[3pt]
ROW & SLACK OR SURPLUS & DUAL PRICES\\
2) & 2.000000 & 0.000000\\
3) & 0.000000 & 2.000000\\
4) & 0.000000 & $-$1.000000\\
5) & 5.000000 & 0.000000\\
\end{tabular}
\end{center}
The signs follow the rule of Step~2: $+2$ on the $\le$ row~3, $-1$ on the
$\ge$ row~4.

\wstep{Step 8. Certify optimality.}
$x^{*}=(0,1,5)$ satisfies every primal row, $u^{*}=(0,-2,1,0)$ every dual
row and sign condition, and
\[
  W=-3\cdot0+12\cdot(-2)-3\cdot1+10\cdot0=-27=C^{*} .
\]
By weak duality no feasible plan costs less than $-27$, so both solutions
are optimal.
\wtrap{The certificate needs the dual of the LP \emph{as stated}. Taking
the dual of the converted max form instead is fine too, but then all its
variables are $\ge 0$ and equal the $z$-row entries of Step~10; mixing the
two conventions is how sign errors creep in.}

\wstep{Step 9. Interpret two numbers.}
Reduced cost $6$ of X1: the cost coefficient $4$ would have to fall by more
than $6$, below $-2$, before $x_{1}$ enters the plan; equivalently, forcing
one unit of $x_{1}$ into the solution raises the minimal cost by $6$. Dual
price $-1$ of row~4: raising its right-hand side from $-3$ to $-2$ raises
the minimal cost by $1$, valid for $-6\le b_{4}\le 3$.

\wstep{Step 10. Rebuild the end tableau.}
In the max form of Step~1 with slacks $s_{1},\dots,s_{4}$ for rows 2--5,
the basis consists of the positive variables $x_{2},x_{3}$ and the slacks of
the nonbinding rows, $s_{1}=2$ and $s_{4}=5$. The $z$-row carries the
reduced costs under the structural columns and, under $s_{i}$, the printed
dual price for a $\le$ row or its negative for a $\ge$ row: $0$, $2$, $1$,
$0$. Bottom right: $-C^{*}=27$. A dot marks an entry the report does not
contain:
\begin{center}\footnotesize
\begin{tabular}{c|ccccccc|c}
Basis & $x_{1}$ & $x_{2}$ & $x_{3}$ & $s_{1}$ & $s_{2}$ & $s_{3}$ & $s_{4}$ & RHS\\
\hline
$x_{2}$ & $\cdot$ & $1$ & $0$ & $0$ & $\cdot$ & $\cdot$ & $0$ & $1$\\
$x_{3}$ & $\cdot$ & $0$ & $1$ & $0$ & $\cdot$ & $\cdot$ & $0$ & $5$\\
$s_{1}$ & $\cdot$ & $0$ & $0$ & $1$ & $\cdot$ & $\cdot$ & $0$ & $2$\\
$s_{4}$ & $\cdot$ & $0$ & $0$ & $0$ & $\cdot$ & $\cdot$ & $1$ & $5$\\
\hline
$z$ & $6$ & $0$ & $0$ & $0$ & $2$ & $1$ & $0$ & $27$\\
\end{tabular}
\end{center}
Check: $z^{*}=\sum_{i}(\text{$z$-row under }s_{i})\cdot b_{i}$ with the
$\le$-form right-hand sides $(3,12,3,10)$ gives
$0\cdot3+2\cdot12+1\cdot3+0\cdot10=27$. \checkmark
\end{worked}

\section{Practice Problems}

Three data sets are shared by the problems below.

\medskip
\noindent\textbf{Workshop LP (Problems 1--6 and 16).} A workshop makes two
products with margins $4$ and $3$ per unit:
\[
  \max\; z=4x_{1}+3x_{2}
  \quad\text{s.t.}\quad
  \begin{array}{rcl}
    2x_{1}+x_{2} &\le& 10 \quad\text{(machine hours)}\\
    x_{1}+x_{2}  &\le& 7  \quad\text{(assembly hours)}\\
    x_{1},x_{2}  &\ge& 0 .
  \end{array}
\]

\medskip
\noindent\textbf{Three-product LP (Problems 7--9).} For
\[
  \max\; z=5x_{1}+4x_{2}+3x_{3}
  \quad\text{s.t.}\quad
  \begin{array}{rcl}
    2x_{1}+3x_{2}+x_{3} &\le& 5\\
    4x_{1}+x_{2}+2x_{3} &\le& 11\\
    3x_{1}+4x_{2}+2x_{3} &\le& 8\\
    x_{1},x_{2},x_{3} &\ge& 0
  \end{array}
\]
with slack variables $s_{1},s_{2},s_{3}$, the simplex method terminates
with the final dictionary
\[
  \begin{array}{rcl}
    z     &=& 13-3x_{2}-s_{1}-s_{3}\\
    x_{1} &=& 2-2x_{2}-2s_{1}+s_{3}\\
    s_{2} &=& 1+5x_{2}+2s_{1}\\
    x_{3} &=& 1+x_{2}+3s_{1}-2s_{3} .
  \end{array}
\]

\medskip
\noindent\textbf{Solver report (Problems 10--15).} A firm makes products
A and B with margins $30$ and $20$:
\[
  \max\; z=30x_{A}+20x_{B}
  \quad\text{s.t.}\quad
  \begin{array}{rcl}
    2x_{A}+x_{B} &\le& 100 \quad\text{(machine hours)}\\
    x_{A}+x_{B}  &\le& 80  \quad\text{(labour hours)}\\
    x_{A}        &\le& 40  \quad\text{(demand for A)}\\
    x_{A},x_{B}  &\ge& 0 .
  \end{array}
\]
The solver prints the following sensitivity report.

\begin{center}
\begin{tabular}{|l|r|r|r|r|r|}
\hline
Variable & Final & Reduced & Objective & Allowable & Allowable\\
         & value & cost    & coeff.    & increase  & decrease\\
\hline
$x_{A}$ & 20 & 0 & 30 & 10 & 10\\
$x_{B}$ & 60 & 0 & 20 & 10 & 5\\
\hline
\end{tabular}

\medskip
\begin{tabular}{|l|r|r|r|r|r|}
\hline
Constraint & Final & Shadow & RHS & Allowable & Allowable\\
           & value & price  &     & increase  & decrease\\
\hline
Machine  & 100 & 10 & 100 & 20 & 20\\
Labour   & 80  & 10 & 80  & 20 & 20\\
Demand A & 20  & 0  & 40  & \texttt{1E+30} & 20\\
\hline
\end{tabular}
\end{center}

\begin{problems}
  \item Solve the workshop LP by corner enumeration. State the optimal
        plan, the optimal value, and which constraints are binding.
  \item Determine both shadow prices of the workshop LP twice: (a) from
        the dual equations of the binding constraints, (b) by raising each
        capacity by one unit and re-solving. Interpret the numbers.
  \item RHS ranging for the workshop LP: for each capacity, express the
        optimal corner as a function of that capacity, derive the validity
        range of its shadow price (allowable increase and decrease), and
        give $z^{*}$ at both ends of each range.
  \item Objective ranging for the workshop LP: over which intervals may
        $c_{1}$ (currently $4$) and $c_{2}$ (currently $3$) vary, one at a
        time, without moving the optimal corner? What happens at the
        endpoints of the ranges?
  \item A neighbouring firm offers the workshop extra hours (all input
        costs are already counted in the margins): machine hours at a
        premium of $1.5$ per hour, or assembly hours at a premium of
        $1.5$ per hour. Which offer should the workshop accept, for at
        most how many hours, and with what total net gain?
  \item A third product is proposed for the workshop: each unit needs
        $1$ machine hour and $2$ assembly hours. (a)~Would it enter the
        plan at a margin of $6$? (b)~At a margin of $4.5$? (c)~At which
        margin is it exactly break-even?
  \item From the final dictionary of the three-product LP read off: the
        optimal plan, the optimal value, the slacks, the three shadow
        prices, and the reduced cost of every variable.
  \item Three-product LP: (a)~By how much may the margin of product~2
        (currently $c_{2}=4$) rise before the plan changes, and what
        happens beyond that point? (b)~Over which range of $b_{2}$
        (currently $11$) is the second shadow price valid, and how do
        plan and value react inside that range?
  \item Three-product LP: (a)~Derive the validity range of the third
        shadow price by perturbing $b_{3}=8$ by $\Delta$ in the final
        dictionary. (b)~For $b_{3}=9$ give the new optimal plan and value
        without re-solving from scratch, and verify them in the original
        constraints.
  \item From the solver report read off: the optimal plan and profit,
        which constraints are binding, the slack of each nonbinding
        constraint, and the meaning of the entry \texttt{1E+30}.
  \item Overtime can be booked at a premium of $8$ per labour hour, up to
        $30$ hours (wages for regular hours are already counted in the
        margins). How many hours should the firm book, and what is the
        resulting net gain? Why not all $30$?
  \item A new product C with margin $25$ would use $1$ machine hour and
        $2$ labour hours per unit. Should it be launched? Give its reduced
        cost and the break-even margin.
  \item Simultaneous capacity changes: (a)~machine $+10$ hours and labour
        $+10$ hours; (b)~machine $+15$ hours and labour $-15$ hours. For
        each case: does the 100\%-rule certify the report's shadow prices,
        and if so, what is the new optimal profit?
  \item Margin changes, one at a time: (a)~the margin of B rises from
        $20$ to $28$: new plan and profit? (b)~the margin of A falls
        from $30$ to $18$: what does the report tell you, and what is
        the new optimum (found by corner enumeration)?
  \item (Beyond the range, solver report.) (a)~Machine hours rise from
        $100$ to $130$, although the allowable increase is only $20$.
        What does the report guarantee for the first $20$ hours, and which
        lower and upper bound does it give for $z^{*}$ at $130$ hours?
        Find the true optimum by corner enumeration. (b)~Machine hours fall
        from $100$ to $70$ (allowable decrease $20$). Give the exact loss
        of the first $20$ hours and a bound for $z^{*}$ at $70$ hours;
        compare with the true optimum.
  \item (Beyond the range, workshop LP.) Using the shadow prices of
        Problem~2 and the ranges of Problem~3, bound $z^{*}$ without
        re-solving, then compute the true value: (a)~assembly capacity
        falls from $7$ to $3$; (b)~assembly capacity rises from $7$ to
        $12$; (c)~machine capacity rises from $10$ to $16$. (d)~The margin
        $c_{1}$ rises from $4$ to $7$, beyond its range of Problem~4:
        which bound does the old plan give, and what is the new optimum?
\end{problems}

\section{Solutions}

\begin{solutions}
  \item The corners of the feasible region and their values:
        \[
          \begin{array}{c|c}
            (x_{1},x_{2}) & z=4x_{1}+3x_{2}\\\hline
            (0,0) & 0\\
            (5,0) & 20\\
            (3,4) & 24\\
            (0,7) & 21
          \end{array}
        \]
        The corner $(3,4)$ solves $2x_{1}+x_{2}=10$, $x_{1}+x_{2}=7$
        (subtracting gives $x_{1}=3$, hence $x_{2}=4$). Optimal plan
        $x^{*}=(3,4)$ with $z^{*}=24$; \emph{both} constraints are
        binding ($2\cdot 3+4=10$, $3+4=7$).

  \item (a) Both constraints bind and both variables are basic, so the
        duals solve the pricing-out equations
        \[
          2y_{1}+y_{2}=4,\qquad y_{1}+y_{2}=3 .
        \]
        Subtracting: $y_{1}=1$, hence $y_{2}=2$.

        (b) Machine $10\to 11$: the corner becomes $x_{1}=11-7=4$,
        $x_{2}=14-11=3$, $z=16+9=25$, so $\Delta z=1=y_{1}$. Assembly
        $7\to 8$: $x_{1}=10-8=2$, $x_{2}=16-10=6$, $z=8+18=26$, so
        $\Delta z=2=y_{2}$.

        Interpretation: one extra machine hour is worth $1$, one extra
        assembly hour is worth $2$: assembly is the scarcer bottleneck
        per hour.

  \item Keep both constraints binding and re-solve the $2\times2$ system
        with a parametric right-hand side.

        \emph{Machine capacity $b_{1}$} (assembly fixed at $7$): from
        $2x_{1}+x_{2}=b_{1}$, $x_{1}+x_{2}=7$ we get
        \[
          x_{1}=b_{1}-7,\qquad x_{2}=14-b_{1}.
        \]
        Nonnegativity forces $7\le b_{1}\le 14$: allowable decrease $3$,
        allowable increase $4$. On this range
        $z^{*}=4(b_{1}-7)+3(14-b_{1})=b_{1}+14$, slope $y_{1}=1$: at
        $b_{1}=7$, $x^{*}=(0,7)$ and $z^{*}=21$; at $b_{1}=14$,
        $x^{*}=(7,0)$ and $z^{*}=28$.

        \emph{Assembly capacity $b_{2}$} (machine fixed at $10$): from
        $2x_{1}+x_{2}=10$, $x_{1}+x_{2}=b_{2}$ we get
        \[
          x_{1}=10-b_{2},\qquad x_{2}=2b_{2}-10 .
        \]
        Nonnegativity forces $5\le b_{2}\le 10$: allowable decrease $2$,
        allowable increase $3$. On this range
        $z^{*}=4(10-b_{2})+3(2b_{2}-10)=2b_{2}+10$, slope $y_{2}=2$: at
        $b_{2}=5$, $x^{*}=(5,0)$ and $z^{*}=20$; at $b_{2}=10$,
        $x^{*}=(0,10)$ and $z^{*}=30$.

  \item The binding constraints have slopes $-2$ (machine) and $-1$
        (assembly); the corner $(3,4)$ stays optimal while
        \[
          -2\;\le\;-\frac{c_{1}}{c_{2}}\;\le\;-1
          \qquad\Longleftrightarrow\qquad
          1\;\le\;\frac{c_{1}}{c_{2}}\;\le\;2 .
        \]
        With $c_{2}=3$ fixed: $3\le c_{1}\le 6$ (allowable decrease $1$,
        increase $2$). With $c_{1}=4$ fixed: $2\le c_{2}\le 4$ (allowable
        decrease $1$, increase $1$). At an endpoint the iso-profit line is
        parallel to a binding constraint and a whole edge is optimal:
        e.g.\ $c_{1}=3$ ties $(3,4)$ with $(0,7)$ (both give $z=21$), and
        $c_{1}=6$ ties $(3,4)$ with $(5,0)$ (both give $z=30$). Beyond an
        endpoint the optimum jumps to that neighbouring corner.

  \item Compare each premium with the shadow price. Machine hours:
        premium $1.5>y_{1}=1$: decline; each bought hour would lose
        $0.5$. Assembly hours: premium $1.5<y_{2}=2$, accept; each
        hour nets $2-1.5=0.5$. The shadow price is guaranteed only up to
        the allowable increase of $b_{2}$, which is $3$ (Problem~3), so
        buy $3$ assembly hours for a net gain of $3\cdot 0.5=1.5$. At
        $b_{2}=10$ the plan is $(0,10)$ and assembly stops being worth
        $2$: for $b_{2}>10$ the optimum stays at $(0,10)$, assembly hours
        go slack and their shadow price drops to $0$, so further hours at
        a positive premium are a pure loss.

  \item Price out the new product against the duals $y=(1,2)$:
        opportunity cost of one unit $=1\cdot y_{1}+2\cdot y_{2}
        =1+4=5$.

        (a) Margin $6>5$: reduced cost $6-5=+1>0$, so the product
        \emph{would} enter the plan (the current plan is no longer
        optimal).

        (b) Margin $4.5<5$: reduced cost $-0.5$, the product stays out
        and the plan $(3,4)$ remains optimal.

        (c) Break-even at margin $=5$: entering neither helps nor hurts
        (multiple optima).

  \item Setting the nonbasic variables $x_{2}=s_{1}=s_{3}=0$ in the
        dictionary:
        \[
          x^{*}=(2,\,0,\,1),\qquad z^{*}=13,\qquad
          s^{*}=(0,\,1,\,0).
        \]
        Constraints $1$ and $3$ bind, constraint $2$ has slack $1$. The
        shadow prices are the negatives of the slack coefficients in the
        $z$-row: $y=(1,\,0,\,1)$. Reduced costs: $r_{1}=r_{3}=0$ (basic),
        $r_{2}=-3$ (the coefficient of $x_{2}$ in the $z$-row). Check via
        pricing out: $r_{2}=4-(3\cdot1+1\cdot0+4\cdot1)=4-7=-3$.

  \item (a) The reduced cost of $x_{2}$ is $-3$, so $c_{2}$ may rise by
        $3$, from $4$ up to $7$, before anything changes. At $c_{2}=7$
        product~2 is exactly break-even (alternative optima); for
        $c_{2}>7$ the variable $x_{2}$ enters the basis and the plan
        changes. Any \emph{decrease} of $c_{2}$ is harmless: $x_{2}$
        simply stays at $0$.

        (b) Constraint $2$ is nonbinding with slack $s_{2}=1$ and
        $y_{2}=0$. Its right-hand side may decrease by the slack, $1$
        (down to $10$), or increase without limit:
        allowable decrease $1$, allowable increase $\infty$. Inside this
        range neither the plan $(2,0,1)$ nor the value $13$ moves,
        only the slack $s_{2}=b_{2}-10$ changes.

  \item (a) Raising $b_{3}$ by $\Delta$ has the same effect on the basic
        variables as setting $s_{3}=-\Delta$ in the dictionary:
        \[
          x_{1}=2-\Delta,\qquad s_{2}=1,\qquad x_{3}=1+2\Delta,\qquad
          z=13+\Delta .
        \]
        Feasibility of the basis requires $x_{1}\ge0$ and $x_{3}\ge0$:
        \[
          2-\Delta\ge 0\;\Rightarrow\;\Delta\le 2,\qquad
          1+2\Delta\ge 0\;\Rightarrow\;\Delta\ge -\tfrac12 .
        \]
        So $y_{3}=1$ is valid for $b_{3}\in[7.5,\,10]$: allowable
        decrease $0.5$, allowable increase $2$.

        (b) $b_{3}=9$ means $\Delta=1$ (inside the range):
        \[
          x^{*}=(1,\,0,\,3),\qquad z^{*}=13+1=14 .
        \]
        Check in the original constraints: $2\cdot1+3=5$ (binding),
        $4\cdot1+2\cdot3=10\le 11$ (slack $1$), $3\cdot1+2\cdot3=9$
        (binding); $z=5\cdot1+3\cdot3=14$. All confirmed.

  \item Optimal plan $x_{A}=20$, $x_{B}=60$, profit
        $z^{*}=30\cdot20+20\cdot60=1800$. Machine and labour are binding:
        their final values equal their right-hand sides ($100$ and $80$)
        and their shadow prices are positive. The demand constraint is
        nonbinding: final value $20$ against a limit of $40$, slack
        $20$, shadow price $0$. \texttt{1E+30} is the solver's notation
        for $\infty$: the demand limit can be raised without any effect
        on the optimum, so there is no finite allowable increase.

  \item The premium $8$ is below the labour shadow price $10$, so
        overtime pays: but the shadow price is guaranteed only for the
        allowable increase of $20$ hours. Book $20$ hours: net gain
        $(10-8)\cdot 20=40$. Hours beyond the $20$th fall outside the
        validity range; there the basis changes ($x_{A}$ has been driven
        to $0$) and extra labour becomes worthless (its shadow price
        drops to $0$) while still costing the premium. Hence not all
        $30$.

  \item Opportunity cost of one unit of C at the current shadow prices:
        \[
          1\cdot 10+2\cdot 10=30>25=\text{margin}.
        \]
        Do not launch: the capacity C would consume is worth more in the
        current plan than C's margin. Reduced cost $r_{C}=25-30=-5$;
        break-even margin $30$ (above $30$, C would enter the plan).

  \item (a) Fractions of the allowable changes:
        $\tfrac{10}{20}+\tfrac{10}{20}=1\le 1$, so the 100\%-rule
        certifies the shadow prices. New profit:
        \[
          z^{*}=1800+10\cdot 10+10\cdot 10=2000 .
        \]
        (b) $\tfrac{15}{20}+\tfrac{15}{20}=1.5>1$: the rule is silent and
        the report alone cannot answer. Indeed the basis breaks: keeping
        machine and labour binding would give
        $x_{A}=115-65=50>40$, violating the demand limit. Re-solving
        with $b=(115,65,40)$ by corner enumeration gives the optimum
        $x=(40,25)$ with $z=30\cdot40+20\cdot25=1700$, not the naive
        $1800+10\cdot15-10\cdot15=1800$.

  \item (a) The increase of $8$ (from $20$ to $28$) is within B's
        allowable increase $10$, so the plan $(20,60)$ is unchanged and
        \[
          z^{*}=1800+8\cdot x_{B}^{*}=1800+8\cdot 60=2280 .
        \]
        (b) The decrease of $12$ exceeds A's allowable decrease $10$:
        the report only says the plan changes; it does not say to what.
        Corner enumeration with $c=(18,20)$:
        \[
          \begin{array}{c|c}
            (x_{A},x_{B}) & z=18x_{A}+20x_{B}\\\hline
            (0,0) & 0\\
            (40,0) & 720\\
            (40,20) & 1120\\
            (20,60) & 1560\\
            (0,80) & 1600
          \end{array}
        \]
        New optimum $x=(0,80)$ with $z^{*}=1600$: product A is dropped
        entirely.

  \item (a)~Machine hours have shadow price $10$ on $[80,120]$. The first
        $20$ extra hours are exact: $z^{*}(120)=1800+10\cdot20=2000$.
        Beyond $120$ every further hour is worth \emph{at most} $10$, and
        extra capacity never hurts, so
        \[
          2000\;\le\; z^{*}(130)\;\le\; 2000+10\cdot 10=2100 .
        \]
        True optimum at $b_{1}=130$: the corners of $2x_{A}+x_{B}\le 130$,
        $x_{A}+x_{B}\le 80$, $x_{A}\le 40$ are $(0,0)$, $(40,0)$, $(40,40)$
        and $(0,80)$, worth $0$, $1200$, $2000$, $1600$. So
        $z^{*}(130)=2000$ at $(40,40)$, using only $120$ machine hours: the
        demand limit of A now binds and the last $10$ hours are worthless,
        the bottom end of the bound.

        (b)~The first $20$ hours lost cost exactly $10$ each:
        $z^{*}(80)=1800-200=1600$. Every further hour lost costs
        \emph{at least} $10$, so $z^{*}(70)\le 1600-100=1500$. True
        optimum at $b_{1}=70$: corners $(0,0)$, $(35,0)$, $(0,70)$ (labour,
        $70\le 80$, and demand no longer bind), worth $0$, $1050$, $1400$.
        So $z^{*}(70)=1400\le 1500$: below $80$ hours a machine hour is
        worth $20$ (it is all spent on B, which earns $20$ per machine
        hour), twice the reported price.

  \item The shadow prices are $y=(1,2)$, valid for $7\le b_{1}\le 14$ and
        $5\le b_{2}\le 10$, with $z^{*}=24$.

        (a)~The first $2$ hours of the decrease ($7\to 5$) cost exactly
        $2\cdot2=4$, so $z^{*}(5)=20$; each of the remaining $2$ costs at
        least $2$, so $z^{*}(3)\le 20-4=16$. True: with
        $x_{1}+x_{2}\le 3$ the machine row ($2\cdot3=6\le 10$) is slack and
        the best corner is $(3,0)$ with $z^{*}=12\le 16$.

        (b)~The first $3$ hours ($7\to 10$) gain exactly $2\cdot3=6$:
        $z^{*}(10)=30$. The remaining $2$ hours gain at most $2$ each, so
        $30\le z^{*}(12)\le 34$. True: corners of $2x_{1}+x_{2}\le 10$,
        $x_{1}+x_{2}\le 12$ are $(0,0)$, $(5,0)$, $(0,10)$, worth $0$,
        $20$, $30$; $z^{*}(12)=30$, the lower end: assembly has gone slack.

        (c)~The first $4$ machine hours ($10\to 14$) gain exactly
        $1\cdot4=4$: $z^{*}(14)=28$. The remaining $2$ gain at most $1$
        each: $28\le z^{*}(16)\le 30$. True: with $2x_{1}+x_{2}\le 16$ and
        $x_{1}+x_{2}\le 7$ the corners are $(0,0)$, $(7,0)$, $(0,7)$, worth
        $0$, $28$, $21$; $z^{*}(16)=28$.

        (d)~The range of $c_{1}$ is $[3,6]$, so $7$ lies outside. The old
        plan $(3,4)$ stays feasible and earns $7\cdot3+3\cdot4=33$, hence
        $z^{*}\ge 33$. Corners: $(5,0)$ gives $35$, $(3,4)$ gives $33$,
        $(0,7)$ gives $21$. New optimum $(5,0)$ with $z^{*}=35\ge 33$.
\end{solutions}

\section{Additional Problems}

Problems 1--5 refer to the LP
\[
  \max\; z=2x_{1}+3x_{2}
  \quad\text{s.t.}\quad
  \begin{array}{rcl}
    x_{1}+x_{2}  &\le& 5 \quad\text{(resource 1)}\\
    x_{1}+2x_{2} &\le& 8 \quad\text{(resource 2)}\\
    x_{1},x_{2}  &\ge& 0 .
  \end{array}
\]

\begin{problems}
  \item Solve the LP by corner enumeration; state the binding constraints
        and compute both shadow prices from the dual equations.
  \item Derive the validity range of each shadow price (allowable
        increase and decrease of $b_{1}=5$ and of $b_{2}=8$) and give
        $z^{*}$ at the ends of each range.
  \item Over which intervals may $c_{1}=2$ and $c_{2}=3$ vary, one at a
        time, without moving the optimal corner?
  \item Extra units are offered (input costs already counted in the
        margins): resource 1 at a premium of $0.8$ per unit, resource 2
        at a premium of $1.5$ per unit. Which offer to accept, for how
        many units, and with what total net gain?
  \item A new product with margin $4$ would use $2$ units of resource 1
        and $1$ unit of resource 2 per unit. Would it enter the plan?
        Give its reduced cost and the largest margin at which it stays
        out.
  \item For the LP $\max z=6x_{1}+5x_{2}$ subject to $x_{1}+x_{2}\le 5$
        (constraint $C_{1}$), $3x_{1}+2x_{2}\le 12$ (constraint $C_{2}$),
        $x_{1},x_{2}\ge0$, the solver reports: final values
        $x_{1}=2$, $x_{2}=3$, both reduced costs $0$; shadow prices
        $y_{1}=3$ (RHS $5$, allowable increase $1$, decrease $1$) and
        $y_{2}=1$ (RHS $12$, allowable increase $3$, decrease $2$).
        (a)~Is it worth buying one extra unit of $C_{1}$'s resource at a
        premium of $2$? (b)~Both right-hand sides rise by $1$
        simultaneously: what does the 100\%-rule say, and what actually
        happens (check directly)?
  \item A firm makes three products from machine hours (row~2) and labour
        hours (row~4) and has promised a customer at least $7$ units of
        product~2 (row~3):
        \[
          \max\; Z=6x_{1}+9x_{2}+7x_{3}
          \quad\text{s.t.}\quad
          \begin{array}{rcl}
            2x_{1}+2x_{2}+x_{3} &\le& 20\\
            x_{2} &\ge& 7\\
            2x_{1}+x_{2}+2x_{3} &\le& 25\\
            x &\ge& 0 .
          \end{array}
        \]
        The solver report is incomplete:
        \begin{center}\footnotesize\ttfamily
        \begin{tabular}{lrr}
        \multicolumn{3}{l}{OBJECTIVE FUNCTION VALUE\quad 1)\ \underline{\hspace{3em}}}\\[3pt]
        VARIABLE & VALUE & REDUCED COST\\
        X1 & \underline{\hspace{3em}} & 8.000000\\
        X2 & \underline{\hspace{3em}} & 0.000000\\
        X3 & \underline{\hspace{3em}} & 0.000000\\[3pt]
        ROW & SLACK OR SURPLUS & DUAL PRICES\\
        2) & 0.000000 & \underline{\hspace{3em}}\\
        3) & 0.000000 & \underline{\hspace{3em}}\\
        4) & \underline{\hspace{3em}} & \underline{\hspace{3em}}\\
        \end{tabular}

        \medskip
        \begin{tabular}{lrrr}
        \multicolumn{4}{l}{OBJ COEFFICIENT RANGES}\\
        VARIABLE & CURRENT COEF & ALLOWABLE INCREASE & ALLOWABLE DECREASE\\
        X1 & 6.000000 & 8.000000 & INFINITY\\
        X2 & 9.000000 & 5.000000 & INFINITY\\
        X3 & 7.000000 & INFINITY & 2.500000\\[3pt]
        \multicolumn{4}{l}{RIGHTHAND SIDE RANGES}\\
        ROW & CURRENT RHS & ALLOWABLE INCREASE & ALLOWABLE DECREASE\\
        2 & 20.000000 & 3.000000 & 6.000000\\
        3 & 7.000000 & 3.000000 & 2.000000\\
        4 & 25.000000 & INFINITY & 6.000000\\
        \end{tabular}
        \end{center}
        (a)~Fill in every blank, using the reduced costs, the slacks and
        complementary slackness. (b)~Write the dual and certify
        optimality. (c)~Interpret the dual price of row~3 and the reduced
        cost of X1. (d)~X2 is basic, yet its allowable decrease is
        \texttt{INFINITY}. Why?
  \item The minimum problem
        \[
          \min\; C=-6x_{1}+4x_{2}+5x_{3}
          \quad\text{s.t.}\quad
          \begin{array}{rcll}
            -2x_{1}+x_{2} &\ge& -4 & \text{(row 2)}\\
            x_{1}+2x_{2}+x_{3} &\le& 28 & \text{(row 3)}\\
            2x_{1}+2x_{3} &\ge& 8 & \text{(row 4)}\\
            x &\ge& 0
          \end{array}
        \]
        produced the incomplete report
        \begin{center}\footnotesize\ttfamily
        \begin{tabular}{lrr}
        \multicolumn{3}{l}{OBJECTIVE FUNCTION VALUE\quad 1)\ \underline{\hspace{3em}}}\\[3pt]
        VARIABLE & VALUE & REDUCED COST\\
        X1 & \underline{\hspace{3em}} & 0.000000\\
        X2 & \underline{\hspace{3em}} & 0.000000\\
        X3 & \underline{\hspace{3em}} & 3.000000\\[3pt]
        ROW & SLACK OR SURPLUS & DUAL PRICES\\
        2) & 0.000000 & \underline{\hspace{3em}}\\
        3) & \underline{\hspace{3em}} & \underline{\hspace{3em}}\\
        4) & 0.000000 & \underline{\hspace{3em}}\\
        \end{tabular}

        \medskip
        \begin{tabular}{lrrr}
        \multicolumn{4}{l}{RIGHTHAND SIDE RANGES}\\
        ROW & CURRENT RHS & ALLOWABLE INCREASE & ALLOWABLE DECREASE\\
        2 & $-$4.000000 & 8.000000 & 4.000000\\
        3 & 28.000000 & INFINITY & 16.000000\\
        4 & 8.000000 & 6.400000 & 4.000000\\
        \end{tabular}
        \end{center}
        (a)~Is the LP a standard maximum problem? (b)~Write the dual (sign
        rules for a min problem). (c)~Fill in every blank and certify
        optimality. (d)~What does the value $3$ in the column
        \texttt{REDUCED COST} of X3 mean? (e)~How does the minimal cost
        change if the right-hand side of row~2 rises from $-4$ to $2$, and
        what can be said if that of row~4 rises from $8$ to $18$?
  \item Rebuild the end tableau of Additional Problem~7 from the completed
        report: (a)~as far as the report allows (basis, RHS column,
        $z$-row), marking the entries it does not contain; (b)~then
        complete the body from the model, by expressing the basic
        variables through the nonbasic ones.
  \item Rebuild the end tableau of the solver-report LP of Problems
        10--15 (slacks $s_{1},s_{2},s_{3}$ for machine, labour, demand)
        from its report alone, then complete the body using the end
        tableau of the two-constraint running example (Chapter~4, worked
        example, Step~6) and $s_{3}=40-x_{A}$.
  \item The farm model with hired labour,
        \[
          \max\; Z=300x_{1}+420x_{2}-9h
          \quad\text{s.t.}\quad
          \begin{array}{rcll}
            x_{1}+x_{2} &\le& 60 & \text{(land)}\\
            8x_{1}+12x_{2}-h &\le& 400 & \text{(labour)}\\
            h &\le& 240 & \text{(hiring cap)},
          \end{array}
        \]
        $x_{1},x_{2},h\ge 0$, returns the report
        \begin{center}\footnotesize\ttfamily
        \begin{tabular}{lrr}
        \multicolumn{3}{l}{OBJECTIVE FUNCTION VALUE\quad 1)\ 20640.00}\\[3pt]
        VARIABLE & VALUE & REDUCED COST\\
        X1 & 20.000000 & 0.000000\\
        X2 & 40.000000 & 0.000000\\
        H & 240.000000 & 0.000000\\[3pt]
        ROW & SLACK OR SURPLUS & DUAL PRICES\\
        2) & 0.000000 & 60.000000\\
        3) & 0.000000 & 30.000000\\
        4) & 0.000000 & 21.000000\\
        \end{tabular}

        \medskip
        \begin{tabular}{lrrr}
        \multicolumn{4}{l}{OBJ COEFFICIENT RANGES}\\
        VARIABLE & CURRENT COEF & ALLOWABLE INCREASE & ALLOWABLE DECREASE\\
        X1 & 300.0000 & 84.00000 & 20.00000\\
        X2 & 420.0000 & 30.00000 & 84.00000\\
        H & $-$9.000000 & INFINITY & 21.00000\\[3pt]
        \multicolumn{4}{l}{RIGHTHAND SIDE RANGES}\\
        ROW & CURRENT RHS & ALLOWABLE INCREASE & ALLOWABLE DECREASE\\
        2 & 60.00000 & 20.00000 & 6.666667\\
        3 & 400.0000 & 80.00000 & 160.0000\\
        4 & 240.0000 & 80.00000 & 160.0000\\
        \end{tabular}
        \end{center}
        (a)~Explain why the dual price of the labour row ($30$) exceeds the
        wage ($9$), and show that $30=9+21$. (b)~Verify the three dual
        prices by complementary slackness. (c)~A contractor offers extra
        hours beyond the cap at $25$ per hour. Accept? How many hours are
        covered by the report, and with what net gain? (d)~Up to which
        hourly wage does the farm keep hiring all $240$ hours?
  \item (Beyond the range, minimum problem.) For the LP of the second
        worked example ($C^{*}=-27$, dual prices $2$ for row~3 and $-1$
        for row~4, ranges as printed there) bound $C^{*}$ without
        re-solving, then compare with the true values $-55$ and
        $-\tfrac{90}{7}$ respectively: (a)~the right-hand side of row~3
        rises from $12$ to $30$; (b)~the right-hand side of row~4 rises
        from $-3$ to $6$. (c)~The cost coefficient of $x_{2}$ rises from
        $-2$ to $12$: which bound does the old plan give? (The new minimum
        is $-15$.)
  \item True or false for the classic report layout (justify briefly):
        (a)~a dual price is never negative;
        (b)~in a minimum problem, a dual price of $3$ on a $\le$ row means
        that one more unit of right-hand side lowers the minimal cost by
        $3$;
        (c)~a variable with value $0$ and reduced cost $0$ is certainly
        basic;
        (d)~rewriting $\min C$ as $\max\,(-C)$ changes the signs of all
        dual prices;
        (e)~a $\ge$ row with surplus $4$ has allowable increase $4$ and
        allowable decrease \texttt{INFINITY}.
\end{problems}

\section{Additional Solutions}

\begin{solutions}
  \item Corners and values:
        \[
          \begin{array}{c|c}
            (x_{1},x_{2}) & z=2x_{1}+3x_{2}\\\hline
            (0,0) & 0\\
            (5,0) & 10\\
            (2,3) & 13\\
            (0,4) & 12
          \end{array}
        \]
        The corner $(2,3)$ solves $x_{1}+x_{2}=5$, $x_{1}+2x_{2}=8$
        (subtracting gives $x_{2}=3$, hence $x_{1}=2$). Optimum
        $x^{*}=(2,3)$, $z^{*}=13$, both constraints binding. Dual
        equations (both variables basic):
        \[
          y_{1}+y_{2}=2,\qquad y_{1}+2y_{2}=3
          \quad\Longrightarrow\quad y_{2}=1,\; y_{1}=1 .
        \]

  \item \emph{Resource 1} ($b_{2}=8$ fixed): from $x_{1}+x_{2}=b_{1}$,
        $x_{1}+2x_{2}=8$ we get $x_{2}=8-b_{1}$, $x_{1}=2b_{1}-8$.
        Nonnegativity gives $4\le b_{1}\le 8$: allowable decrease $1$,
        increase $3$. There $z^{*}=2(2b_{1}-8)+3(8-b_{1})=b_{1}+8$,
        slope $y_{1}=1$: at $b_{1}=4$, $x^{*}=(0,4)$, $z^{*}=12$; at
        $b_{1}=8$, $x^{*}=(8,0)$, $z^{*}=16$.

        \emph{Resource 2} ($b_{1}=5$ fixed): from $x_{1}+x_{2}=5$,
        $x_{1}+2x_{2}=b_{2}$ we get $x_{2}=b_{2}-5$, $x_{1}=10-b_{2}$.
        Nonnegativity gives $5\le b_{2}\le 10$: allowable decrease $3$,
        increase $2$. There $z^{*}=2(10-b_{2})+3(b_{2}-5)=b_{2}+5$,
        slope $y_{2}=1$: at $b_{2}=5$, $x^{*}=(5,0)$, $z^{*}=10$; at
        $b_{2}=10$, $x^{*}=(0,5)$, $z^{*}=15$.

  \item Binding slopes: $-1$ (resource 1) and $-\tfrac12$ (resource 2);
        the corner $(2,3)$ stays optimal while
        \[
          -1\;\le\;-\frac{c_{1}}{c_{2}}\;\le\;-\frac12
          \qquad\Longleftrightarrow\qquad
          \frac12\;\le\;\frac{c_{1}}{c_{2}}\;\le\;1 .
        \]
        With $c_{2}=3$ fixed: $1.5\le c_{1}\le 3$ (decrease $0.5$,
        increase $1$). With $c_{1}=2$ fixed: $2\le c_{2}\le 4$ (decrease
        $1$, increase $1$). At each endpoint an edge of alternative
        optima appears.

  \item Resource 1: premium $0.8<y_{1}=1$, accept; each unit nets
        $0.2$, guaranteed up to the allowable increase $3$ (Additional
        Problem~2), so buy $3$ units for a net gain of $3\cdot0.2=0.6$.
        Resource 2: premium $1.5>y_{2}=1$, decline.

  \item Opportunity cost per unit at $y=(1,1)$:
        $2\cdot1+1\cdot1=3$. Margin $4>3$, so the product \emph{would}
        enter the plan; its reduced cost is $4-3=+1$. It stays out as
        long as its margin is at most $3$ (at exactly $3$ it is
        break-even).

  \item (a) Premium $2<y_{1}=3$: yes; one unit is within the allowable
        increase $1$ and nets $3-2=1$.

        (b) Fractions of the allowable increases:
        $\tfrac11+\tfrac13=\tfrac43>1$, so the 100\%-rule is silent.
        Direct check: with $b=(6,13)$, keeping both constraints binding
        gives $x_{1}=13-2\cdot6=1$, $x_{2}=3\cdot6-13=5$; both are
        nonnegative, so the same basis is still optimal and the shadow
        prices \emph{did} remain valid:
        \[
          z^{*}=27+3\cdot1+1\cdot1=31=6\cdot1+5\cdot5 .
        \]
        Lesson: the 100\%-rule is sufficient, not necessary: failing
        it does not prove the basis changes, it only withdraws the
        guarantee.

  \item (a)~\emph{Variables.} X1 has reduced cost $8>0$: $x_{1}=0$. X2 and
        X3 are basic. \emph{Rows.} Rows 2 and 3 bind (slack $0$). Three
        rows mean three basic variables; besides $x_{2}$ and $x_{3}$ only
        the slack of row~4 is left, so row~4 is nonbinding, as its
        one-sided range ($\infty$, $6$) confirms.
        \emph{Primal.} Row~3: $x_{2}=7$; row~2: $2\cdot0+14+x_{3}=20$, so
        $x_{3}=6$. Row~4: $0+7+12=19$, slack $25-19=6$ (matching the
        allowable decrease). $Z^{*}=63+42=105$.

        \emph{Dual.} Sign rules for a max problem: rows 2 and 4 ($\le$)
        give $y_{1},y_{3}\ge 0$, row~3 ($\ge$) gives $y_{2}\le 0$:
        \[
          \min\; W=20y_{1}+7y_{2}+25y_{3}
          \quad\text{s.t.}\quad
          2y_{1}+2y_{3}\ge 6,\qquad 2y_{1}+y_{2}+y_{3}\ge 9,\qquad
          y_{1}+2y_{3}\ge 7 .
        \]
        Row~4 is slack: $y_{3}=0$. $x_{3}>0$: $y_{1}=7$. $x_{2}>0$:
        $14+y_{2}=9$, so $y_{2}=-5$. Check X1: $2\cdot7=14\ge 6$ with
        dual slack $8$, the printed reduced cost. \checkmark\ For a max
        problem the dual prices are $y$: row~2 $7$, row~3 $-5$, row~4
        $0$. Completed report: objective value $105$; values
        $0$, $7$, $6$; slack or surplus $0$, $0$, $6$; dual prices $7$,
        $-5$, $0$.

        (b)~$x^{*}=(0,7,6)$ is primal feasible, $y^{*}=(7,-5,0)$ satisfies
        all dual rows and signs, and $W=140-35+0=105=Z^{*}$: both are
        optimal.

        (c)~Dual price $-5$ of row~3: each additional unit promised to the
        customer lowers the maximal profit by $5$ (valid for a promise
        between $5$ and $10$ units). A unit of product~2 earns $9$ but uses
        $2$ machine hours, which are worth $2\cdot7=14$ in product~3.
        Reduced cost $8$ of X1: its margin would have to rise by more than
        $8$, above $14$, before product~1 is made; at present its two
        machine hours (worth $14$) earn only $6$.

        (d)~$x_{2}$ sits at the contract minimum because the contract
        forces it there, not because it is profitable. A lower margin only
        makes product~2 less attractive, and the plan still produces the
        contracted $7$ units and no more: the basis never changes, hence
        \texttt{INFINITY}. An increase by more than $5$ (to above $14$,
        i.e.\ $7$ per machine hour, like product~3) would make it worth
        producing beyond the contract.

  \item (a)~\textbf{No.} Origin test: row~4 reads $0\ge 8$, false, so the
        origin is infeasible and the problem needs phase one. (Rows 2 and
        3 alone would be harmless: row~2 becomes $2x_{1}-x_{2}\le 4$.)

        (b)~Min problem: the $\ge$ rows 2 and 4 get $u_{1},u_{3}\ge 0$, the
        $\le$ row~3 gets $u_{2}\le 0$; the dual constraints are $\le$:
        \[
          \max\; W=-4u_{1}+28u_{2}+8u_{3}
          \quad\text{s.t.}\quad
          -2u_{1}+u_{2}+2u_{3}\le -6,\qquad u_{1}+2u_{2}\le 4,\qquad
          u_{2}+2u_{3}\le 5 .
        \]
        (c)~X3 has reduced cost $3>0$: $x_{3}=0$. Rows 2 and 4 bind:
        $2x_{1}=8$, so $x_{1}=4$, and $-8+x_{2}=-4$, so $x_{2}=4$. Row~3:
        $4+8=12$, slack $16$ (the allowable decrease). $C^{*}=-24+16=-8$.
        Dual: row~3 slack gives $u_{2}=0$; $x_{2}>0$ gives $u_{1}=4$;
        $x_{1}>0$ gives $-8+0+2u_{3}=-6$, so $u_{3}=1$. Signs fine; the
        third dual constraint $0+2=2\le 5$ has slack $3$, the printed
        reduced cost. \checkmark\ Dual prices $=-u$: row~2 $-4$, row~3
        $0$, row~4 $-1$. Certificate: $W=-16+0+8=-8=C^{*}$.

        (d)~The cost coefficient $5$ of $x_{3}$ would have to fall by more
        than $3$, below $2$, before $x_{3}$ enters the plan; forcing one
        unit of $x_{3}$ in raises the minimal cost by $3$. (A unit of
        $x_{3}$ covers $2$ units of row~4, which are worth $2u_{3}=2$, and
        costs $5$.)

        (e)~Row~2: $-4\to 2$ is an increase of $6\le 8$, inside the range.
        The dual price $-4$ means the cost rises by $4$ per unit:
        $C^{*}=-8+24=16$ (new plan $x=(4,10,0)$ from the binding rows,
        with row~3 at $4+20=24\le 28$). Row~4: $8\to 18$ is an increase of
        $10$, beyond the allowable $6.4$. The first $6.4$ units cost exactly
        $1$ each, giving $-8+6.4=-1.6$; each further unit costs \emph{at
        least} $1$, so $C^{*}(18)\ge -1.6+3.6=2$ (if the LP is still
        feasible). Re-solving gives $C^{*}=\tfrac{35}{4}=8.75$, consistent
        with the bound.

  \item (a)~$\le$ form: $2x_{1}+2x_{2}+x_{3}+s_{1}=20$,
        $-x_{2}+s_{2}=-7$, $2x_{1}+x_{2}+2x_{3}+s_{3}=25$. Basis: $x_{2}=7$,
        $x_{3}=6$ and the slack of row~4, $s_{3}=6$. $z$-row: reduced costs
        $8,0,0$; under $s_{1}$ the dual price $7$ (a $\le$ row); under
        $s_{2}$ the negative of $-5$, i.e.\ $5$ (an original $\ge$ row);
        under $s_{3}$ $0$; bottom right $105$.
        \begin{center}\footnotesize
        \begin{tabular}{c|cccccc|c}
        Basis & $x_{1}$ & $x_{2}$ & $x_{3}$ & $s_{1}$ & $s_{2}$ & $s_{3}$ & RHS\\
        \hline
        $x_{2}$ & $\cdot$ & $1$ & $0$ & $\cdot$ & $\cdot$ & $0$ & $7$\\
        $x_{3}$ & $\cdot$ & $0$ & $1$ & $\cdot$ & $\cdot$ & $0$ & $6$\\
        $s_{3}$ & $\cdot$ & $0$ & $0$ & $\cdot$ & $\cdot$ & $1$ & $6$\\
        \hline
        $z$ & $8$ & $0$ & $0$ & $7$ & $5$ & $0$ & $105$\\
        \end{tabular}
        \end{center}
        (b)~From row~3: $x_{2}=7+s_{2}$. From row~2:
        $x_{3}=20-2x_{1}-2x_{2}-s_{1}=6-2x_{1}-s_{1}-2s_{2}$. From row~4:
        $s_{3}=25-2x_{1}-x_{2}-2x_{3}=6+2x_{1}+2s_{1}+3s_{2}$. Moving the
        nonbasic terms to the left gives the rows of the tableau:
        \begin{center}\footnotesize
        \begin{tabular}{c|cccccc|c}
        Basis & $x_{1}$ & $x_{2}$ & $x_{3}$ & $s_{1}$ & $s_{2}$ & $s_{3}$ & RHS\\
        \hline
        $x_{2}$ & $0$ & $1$ & $0$ & $0$ & $-1$ & $0$ & $7$\\
        $x_{3}$ & $2$ & $0$ & $1$ & $1$ & $2$ & $0$ & $6$\\
        $s_{3}$ & $-2$ & $0$ & $0$ & $-2$ & $-3$ & $1$ & $6$\\
        \hline
        $z$ & $8$ & $0$ & $0$ & $7$ & $5$ & $0$ & $105$\\
        \end{tabular}
        \end{center}
        Check of the $z$-row: substituting into $z=6x_{1}+9x_{2}+7x_{3}$
        gives $z=105-8x_{1}-7s_{1}-5s_{2}$, i.e.\ the entries $8,7,5$.
        \checkmark

  \item From the report: basis $x_{A}=20$, $x_{B}=60$ and the slack of the
        demand row, $s_{3}=40-20=20$. The $z$-row holds the reduced costs
        $0,0$, the shadow prices $10,10,0$ and $z^{*}=1800$. The
        two-constraint end tableau supplies the rows of $x_{A}$ and
        $x_{B}$ (the demand row does not involve $s_{1},s_{2}$ and is
        nonbinding, so they are unchanged). Only its body carries over:
        that example has the same two constraints but the margins $40$
        and $30$, so its $z$-row ($10$, $20$, $2600$) does not apply here,
        while the body $B^{-1}A$ depends on the constraints alone. From
        $s_{3}=40-x_{A}=40-(20-s_{1}+s_{2})=20+s_{1}-s_{2}$ follows the
        third row:
        \begin{center}\footnotesize
        \begin{tabular}{c|ccccc|c}
        Basis & $x_{A}$ & $x_{B}$ & $s_{1}$ & $s_{2}$ & $s_{3}$ & RHS\\
        \hline
        $x_{A}$ & $1$ & $0$ & $1$ & $-1$ & $0$ & $20$\\
        $x_{B}$ & $0$ & $1$ & $-1$ & $2$ & $0$ & $60$\\
        $s_{3}$ & $0$ & $0$ & $-1$ & $1$ & $1$ & $20$\\
        \hline
        $z$ & $0$ & $0$ & $10$ & $10$ & $0$ & $1800$\\
        \end{tabular}
        \end{center}
        From the report alone only the unit columns of $x_{A},x_{B},s_{3}$,
        the RHS column and the $z$-row are known; the six entries under
        $s_{1},s_{2}$ in the body need the model.

  \item (a)~The hiring cap binds: the farm would like to hire more at
        $9$ but cannot. One more \emph{own} hour therefore has the full
        value of labour in production, $30$: it lets the farm move
        $\tfrac14$ hectare from crop~1 to crop~2, which gains
        $\tfrac14(420-300)=30$. Of that value, $9$ is what a hired hour
        costs and $21$ is what the farm would net from one more hour of
        hiring: the dual price of the cap. Hence $30=9+21$.

        (b)~All three variables are basic, so all three dual constraints
        bind:
        \[
          y_{1}+8y_{2}=300,\qquad y_{1}+12y_{2}=420,\qquad -y_{2}+y_{3}=-9 .
        \]
        Subtracting the first from the second, $4y_{2}=120$, so
        $y_{2}=30$, $y_{1}=60$ and $y_{3}=21$. Check:
        $60\cdot60+30\cdot400+21\cdot240=3600+12000+5040=20640=Z^{*}$.
        \checkmark

        (c)~Accept: an extra hour raises $Z^{*}$ by $30$ and costs $25$,
        net $5$. The labour row's dual price holds for an increase of up to
        $80$ hours, so the report covers $80$ hours with a net gain of
        $80\cdot5=400$.

        (d)~The objective coefficient of $h$ is $-9$ with allowable
        decrease $21$: the plan, including all $240$ hired hours, stays
        optimal for a wage up to $9+21=30$, the value of an hour of labour.

  \item (a)~Row~3 ($\le$, dual price $2$, allowable increase $12$): the
        first $12$ units are exact, $C^{*}(24)=-27-2\cdot12=-51$. Each of
        the remaining $6$ units lowers the cost by \emph{at most} $2$, and
        more capacity never raises it, so
        \[
          -63=-51-12\;\le\;C^{*}(30)\;\le\;-51 .
        \]
        The true value $-55$ lies inside.

        (b)~Row~4 ($\ge$, dual price $-1$, allowable increase $6$): the
        first $6$ units are exact, $C^{*}(3)=-27+6=-21$. Each of the
        remaining $3$ units raises the cost by \emph{at least} $1$, so
        $C^{*}(6)\ge -21+3=-18$. The true value
        $-\tfrac{90}{7}\approx-12.86$ satisfies the bound.

        (c)~The increase of $14$ exceeds the allowable increase $12$, so
        the plan may change. The old plan $(0,1,5)$ stays feasible and now
        costs $0+12-25=-13$, so $C^{*}\le -13$. The true minimum $-15$
        respects the bound.

  \item (a)~\textbf{False.} A binding $\ge$ row has a dual price
        $\le 0$, as row~4 ($-1$) of the second worked example shows.

        (b)~\textbf{True.} The dual price is the improvement per unit
        increase of the right-hand side, and for a min problem improvement
        means lower cost.

        (c)~\textbf{False.} A nonbasic variable can have reduced cost $0$:
        it can then be brought in without changing the objective value
        (alternative optima). Value $0$ with reduced cost $0$ can also be a
        degenerate basic variable; the report does not tell which.

        (d)~\textbf{False.} Improvement is the same whichever way the
        objective is written; dual prices and reduced costs are unchanged,
        and only the objective value changes sign.

        (e)~\textbf{True.} Raising the requirement by up to the surplus $4$
        keeps the row nonbinding and the basis unchanged; lowering it only
        loosens a row that already has room, without limit.
\end{solutions}
